\documentclass[runningheads]{llncs}

% --------------------------------------------------
% Packages
% --------------------------------------------------
\usepackage{amsmath,amssymb,amsfonts}
\usepackage{graphicx}
\usepackage{booktabs}
\usepackage{array}
\usepackage{multirow}
\usepackage{tabularx}
\usepackage{siunitx}
\usepackage{url}
\usepackage{xurl}
\usepackage{float}
\usepackage{hyperref}

\urlstyle{same}

\title{AI-Driven Forensic Indexing:\\
A Behavioural Model for Chain Snatching Identification}
\titlerunning{AI-Driven Forensic Indexing}
\author{Anonymous Author(s)}
\authorrunning{Anonymous}

\begin{document}

\maketitle

\begin{abstract}
Automated analysis of CCTV footage for chain-snatching incidents is
challenging due to large data volumes, varying visual conditions, and
complex pedestrian--suspect interactions. Existing approaches often
address object detection, motion analysis, depth estimation, or action
recognition independently, and even recent multi-task surveillance
systems rarely jointly address forensic evidence preservation,
computational efficiency, and domain-specific snatch signatures. This
paper presents an AI-driven forensic search framework that integrates
these capabilities into a progressive three-tier architecture,
combining low-level motion and visual preprocessing, semantic object
detection and spatial reasoning, and higher-level behavioural analysis
to identify chain-snatching-related interactions and generate
searchable forensic events, with particular emphasis on modelling
suspect--victim interactions, including pillion-rider behaviour, while
preserving evidence throughout the pipeline. Evaluated on 43 CCTV sequences (35 positive, 8 negative), the complete
system achieved \textbf{100.00\% precision, 94.29\% recall, 97.06\%
F1-score, and 95.35\% accuracy}, while progressive processing reduced
the retained search space to \textbf{2.38\%} of original frames and
preserved \textbf{33 of 35 ground-truth events}. The forensic indexing
component additionally achieved \textbf{100\% search accuracy and
completeness} with sub-millisecond average query latency. These results
demonstrate that progressive visual and behavioural analysis can
substantially reduce CCTV investigation workload while retaining
relevant forensic evidence and enabling efficient event retrieval.
\end{abstract}

\keywords{Chain Snatching Detection, Forensic Surveillance, Computer Vision,
Object Detection, Monocular Depth Estimation, Behavioral Recognition,
Forensic Indexing}

\section{Introduction}

The increasing deployment of CCTV systems has produced large volumes of
surveillance footage that are difficult to inspect manually. In
chain-snatching investigations, relevant events may occur within only a
few frames and involve complex interactions between victims, suspects,
and motorcycles. Frame-level object detection alone is therefore
insufficient for reliable forensic analysis.

Existing approaches often focus on individual surveillance tasks, such as
motion detection, object detection, depth estimation, tracking, or action
recognition~\cite{roy2018snatch}, while some recent systems integrate
multiple vision capabilities~\cite{park2026foresea}. However, these
approaches are not specifically designed for chain-snatching or
domain-specific snatch signatures, and generic action recognition does
not explicitly model the spatial and temporal relationships involved.

This paper presents an AI-driven forensic search framework that
progressively transforms CCTV footage into structured forensic events.
It integrates motion filtering, object detection, tracking, spatial
reasoning, behavioural analysis, evidence fusion, crime-specific
signature analysis, and forensic indexing to reduce the forensic search
space while preserving relevant evidence.

The framework is evaluated on annotated chain-snatching and negative
CCTV sequences. The results demonstrate that the proposed approach can
substantially narrow the forensic candidate space while retaining
relevant crime evidence and supporting efficient CCTV-based
investigation.

\section{Taxonomy of the Literature Survey}

To provide a structured overview of the reviewed literature, the surveyed techniques are categorized into three major research components: preprocessing, localization, and behavioral modeling. Each component is further divided into its respective subdomains based on the primary objective of the methods. Figure~\ref{fig:taxonomy} illustrates the taxonomy adopted in this survey.

\begin{figure}[!htbp]
    \centering
    \includegraphics[width=0.48\textwidth]{chain_snatching_taxonomy.png}
    \caption{Taxonomy of the literature surveyed for AI-driven chain-snatching detection.}
    \label{fig:taxonomy}
\end{figure}

\begin{table}[!t]
\centering
\scriptsize
\renewcommand{\arraystretch}{1.2}
\caption{Comparison of representative low-light image enhancement techniques for forensic surveillance.}
\label{tab:lowlight_comparison}

\begin{tabular}{|l|p{3.0cm}|p{3.5cm}|p{3.8cm}|p{2.5cm}|}
\hline
\textbf{Model} &
\textbf{Core Technique} &
\textbf{Strength} &
\textbf{Primary Limitation} &
\textbf{Suitability} \\ \hline

LIME~\cite{lime} &
Illumination Map Estimation &
Simple illumination correction &
High computational cost; prone to over-exposure under strong lighting &
Limited \\ \hline

Retinex-Net~\cite{retinexnet} &
Deep Retinex Decomposition &
Improves contrast and detail restoration &
Amplifies sensor noise under challenging lighting conditions &
Moderate \\ \hline

EnlightenGAN~\cite{enlightengan} &
GAN-based Image Enhancement &
Produces visually appealing enhanced images &
May generate hallucinated textures, compromising evidentiary reliability &
Low \\ \hline

\textbf{Zero-DCE}~\cite{zerodce} &
Zero-Reference Curve Estimation &
Non-generative, lightweight, and preserves structural details &
Performance may degrade under extremely severe illumination conditions &
High \\ \hline

CycleGAN~\cite{cyclegan} &
GAN-based Image Translation &
Effective style transfer without paired training &
May alter image content through hallucinated structures &
Low \\ \hline

\end{tabular}
\end{table}

\subsection{Low-Light Image Enhancement}
Nighttime CCTV footage suffers from reduced visibility, sensor noise,
and low contrast that obscure hand movements and suspect--victim
interactions, compounded by benchmark datasets such as UCF-Crime \cite{sultani2018ucfcrime} being
predominantly daytime, limiting model robustness under low-light
conditions. Classical illumination-map techniques such as LIME improve
luminance but introduce over-exposure under strong light sources such
as street lamps or headlights~\cite{lime}, while generative
image-to-image approaches such as EnlightenGAN improve perceptual
quality but may hallucinate textures, undesirable in forensic
applications where preserving original visual evidence is essential
~\cite{enlightengan}.

Zero-Reference Deep Curve Estimation (Zero-DCE) instead formulates
enhancement as an image-specific curve-estimation problem, using a
lightweight DCE-Net to predict pixel-wise enhancement curves applied
iteratively:
\begin{equation}
L_{n}(x)=L_{n-1}(x)+\alpha_{n}L_{n-1}(x)
\left(1-L_{n-1}(x)\right),
\label{eq:zerodce}
\end{equation}
where $L_n(x)$ is the enhanced pixel intensity after the
$n^{\text{th}}$ iteration and $\alpha_n$ controls enhancement
magnitude. Trained using non-reference losses (spatial consistency,
exposure control, color constancy, illumination smoothness), Zero-DCE
requires no paired or unpaired reference images~\cite{zerodce}, and its
lightweight, non-generative formulation avoids synthesizing new
detail—making it the most suitable choice for forensic surveillance
among the compared methods.


\subsection{Background Subtraction \& Temporal Filtering}

Background subtraction isolates foreground objects such as pedestrians, suspects, and vehicles from irrelevant background before detection and behavioral analysis, but outdoor surveillance introduces illumination change, camera motion, shadows, and weather noise that reduce localization accuracy. Frame differencing is computationally efficient but sensitive to camera motion and illumination change, often producing fragmented masks~\cite{framediff}. Standard Gaussian Mixture Models (GMM) adapt to gradual background change but lose accuracy under sudden illumination shifts or complex outdoor scenes~\cite{gmm}. Sample-based methods such as ViBe adapt efficiently to changing backgrounds but risk absorbing stationary foreground objects into the background over time~\cite{vibe}. MOG2 extends GMM by adaptively estimating the number of Gaussian components per pixel, improving robustness to gradual illumination change, shadows, and repetitive background motion while retaining real-time performance, and is therefore adopted for motion triage~\cite{zivkovic2004mog2}.


\subsection{Real-Time Object Detection}

Object detection supplies the semantic information required for subsequent depth estimation and spatio-temporal analysis, and must balance accuracy against computational efficiency under motion blur, occlusion, varying scale, and crowding~\cite{malhotra2020objectcomparison}. Two-stage detectors such as R-CNN and Faster R-CNN localize accurately but their computational cost limits real-time use on high-resolution surveillance video~\cite{fasterrcnn}. Single-stage detectors such as SSD and the YOLO family instead perform detection in a single pass for faster inference~\cite{ssd,yolo}; SSD is efficient but can underperform on small or visually challenging objects, while successive YOLO variants progressively improve the accuracy--speed trade-off~\cite{yolov8,yolov10,yolov11}. YOLOv11 in particular incorporates the C3k2 block and C2PSA attention module, improving feature extraction and multi-scale attention while maintaining efficiency, giving it the best accuracy, parameter-efficiency, and inference-speed trade-off for detecting pedestrians and two-wheeled vehicles in challenging footage~\cite{yolov11}.


\subsection{Monocular Depth Estimation}

A major challenge in CCTV-based crime analytics is 2D occlusion, where overlapping bounding boxes may incorrectly suggest physical interaction -- for example, a pedestrian walking near a motorcycle may appear spatially close in image coordinates despite being physically separated, leading to erroneous interaction reasoning and false positives. Stereo configurations and active sensors such as LiDAR or Time-of-Flight (ToF) provide accurate depth but require additional hardware, calibration, and deployment cost, limiting their suitability for large-scale urban surveillance. Deep monocular approaches remove this hardware dependency: MiDaS demonstrated strong cross-dataset relative-depth estimation, but its computational cost and reduced boundary precision in complex scenes limit real-time suitability~\cite{monovit}. Depth-Anything V2~\cite{yang2024depth,liuhe2024depthanything}, which uses a Vision Transformer foundation model trained on large-scale synthetic and real-world data, instead provides robust relative-depth estimation with sharper object boundaries and stronger cross-environment generalization, making it the preferred choice for estimating spatial relationships in complex urban surveillance scenes.


\subsection{Pose Estimation and Tier Integration}

Human pose estimation provides an intermediate skeletal representation that bridges low-level object/depth analysis and higher-level behavioral recognition, and is less sensitive to background clutter and illumination than appearance-based features. MediaPipe Pose is adopted for this role: its heatmap-based keypoint estimation retains useful information under the partial occlusion common during rapid chain-snatching interactions, and its built-in temporal tracking avoids repeated full-scale pose estimation across consecutive frames, reducing computational cost while maintaining short-term skeletal consistency~\cite{mediapipe}. In addition to 2D joint locations, its depth estimates under a weak-perspective camera model are combined with the global depth estimates from Depth-Anything V2 to form a more informative 3D skeleton for the ST-GCN model~\cite{yang2024depth,yan2018liuhe}.


\subsection{Spatio-Temporal Behavioral Recognition}

\begin{table}[!t]
\centering
\scriptsize
\renewcommand{\arraystretch}{1.2}
\caption{Quantitative comparison of representative behavior recognition models for forensic surveillance.}
\label{tab:behavior_quantitative}

\begin{tabular}{|l|c|c|c|}
\hline
\textbf{Metric} &
\textbf{ViViT~\cite{liu2022video}} &
\textbf{2s-AGCN~\cite{shi2019twostreamagcn,wu20222s}} &
\textbf{ST-GCN~\cite{yan2018stgcn}} \\ \hline

Top-1 Accuracy (NTU-120) (\%) $\uparrow$ &
89.1 &
\textbf{91.2} &
89.9 \\ \hline

GFLOPs $\downarrow$ &
50.0 &
18.4 &
\textbf{16.2} \\ \hline

Parameters (Millions) $\downarrow$ &
100+ &
6.9 &
\textbf{3.1} \\ \hline

Skeleton-Based Representation &
No &
Yes &
Yes \\ \hline

Real-Time Deployment &
Moderate &
High &
\textbf{High} \\ \hline

Interpretability &
Low &
High &
\textbf{High} \\ \hline

\end{tabular}
\end{table}

Distinguishing legitimate pedestrian--vehicle interactions from
visually similar chain-snatching events requires modelling body
configuration and motion over time, since object detection and depth
estimation alone cannot determine whether an interaction is criminal.
CNN-LSTM approaches combine spatial and temporal modelling but remain
sensitive to background appearance, illumination, and viewpoint change
~\cite{cnnlstm}; two-stream architectures explicitly model motion via
optical flow, though dense flow computation is costly for continuous
surveillance~\cite{twostreamlstm}; Vision Transformer models such as
ViViT capture long-range temporal dependencies via self-attention but
require comparatively large computational and memory resources
~\cite{liu2022video}.

Graph-based skeleton methods represent the human body as a graph of
joints and anatomical edges: ST-GCN learns spatial and temporal motion
patterns directly from skeletal data, while adaptive variants such as
2s-AGCN improve accuracy at additional computational cost
(Table~\ref{tab:behavior_quantitative})
~\cite{shi2019twostreamagcn,yan2018stgcn}. ST-GCN therefore offers a
balance of cost, accuracy, and deployability, and combined with
MediaPipe Pose and Depth-Anything V2 enables depth-aware
spatio-temporal analysis of pedestrian--motorcycle interactions
~\cite{yan2018stgcn}.

\section{Research Gaps}
\label{sec:gaps}
The comparative analysis in Section~3 identifies several limitations
that hinder the application of existing surveillance approaches to
automated chain-snatching detection. Although substantial progress has
been made in behavioural recognition, existing solutions do not
adequately address the requirements of complex CCTV environments.

\begin{itemize}
\item \textbf{Pillion-Rider Behavioral Gap:} Existing action recognition
approaches primarily target generic human behaviours such as walking,
running, fighting, and falling. Although graph-based approaches such as
ST-GCN effectively model human motion patterns~\cite{yan2018stgcn},
fine-grained interactions involving pedestrians, motorcycles, and rapid
``grab-and-accelerate'' motion in chain-snatching scenarios remain
insufficiently modelled~\cite{roy2018snatch}.

\item \textbf{Nighttime Reliability Gap:} Existing surveillance datasets
are predominantly biased toward daytime conditions, while nighttime
footage suffers from poor illumination, motion blur, and reduced sensor
sensitivity. Traditional enhancement methods may struggle under these
conditions, whereas generative approaches can introduce artifacts not
present in the original footage~\cite{lime,enlightengan,zerodce,retinexnet}.

\item \textbf{Edge Resource Bottleneck:} Increasing recognition
capability is often accompanied by higher computational and memory
requirements, limiting deployment of intensive video-recognition
architectures on resource-constrained edge platforms
~\cite{liu2022video,feichtenhofer2019slowfast}. This motivates
progressive, resource-efficient architectures for forensic video
analysis that reduce unnecessary computation without relying entirely
on cloud-based processing.

\item \textbf{2D Perspective Ambiguity:} Reasoning about interactions
using only two-dimensional image coordinates can produce incorrect
proximity relationships, since objects at different physical depths may
overlap or appear close in the image plane, leading to false
interaction events~\cite{monovit,yang2024depth}. Incorporating monocular
depth estimation can therefore provide additional spatial information
for validating interactions in surveillance footage.
\end{itemize}

\section{Proposed Methodology}
\label{sec:methodology}

\subsection{System Overview}

The proposed approach is three-stage hierarchical in nature where deep learning models of increasing complexity are applied sequentially only when needed. Firstly, motion filtering is applied to discard stationary frames, followed by spatial localization and finally performing spatio-temporal behavioural analysis. This approach reduces the overall computational cost while retaining the forensic accuracy.

\subsection{Motion Triage and Illumination Normalization}
\subsubsection{Background Subtraction using MOG2}

Motion triage uses the MOG2 (Mixture of Gaussians) background subtractor \cite{zivkovic2004mog2}, which models each pixel as an adaptive mixture of Gaussians to separate foreground from background under changing scene conditions. Its adaptive learning rate accommodates gradual illumination changes from weather, while OpenCV's built-in shadow detection masks shadowed regions before further processing.

MOG2 is configured with a history length of 500 frames, a variance threshold of 16.0, and shadow detection enabled. For each incoming frame, a foreground mask is generated and the count of non-zero foreground pixels computed; frames exceeding 5,000 foreground pixels are forwarded to the object detection stage, reducing unnecessary downstream computation.

\subsubsection{Low-Light Enhancement using Zero-DCE}
Frames containing foreground activity are conditionally passed to the
Zero-DCE enhancement network \cite{zerodce} when poor illumination is
detected. The illumination level is estimated as the mean brightness
$\bar{L}(f)=\operatorname{mean}\left(V_{\mathrm{HSV}}(f)\right)$ of the
motion-triaged region in the HSV colour space, and enhancement is applied
only when $\bar{L}(f)<\tau_{lux}$, a tunable threshold selected according
to the deployment environment. Well-illuminated frames are thus retained
without enhancement, reducing unnecessary computation.

Zero-DCE performs enhancement by estimating pixel-wise adjustment curves
rather than directly generating a new image:
\begin{equation}
L_n(x)=L_{n-1}(x)+\alpha_nL_{n-1}(x)
\left(1-L_{n-1}(x)\right),
\end{equation}
where $L_n(x)$ represents the enhanced pixel intensity after the
$n^{\text{th}}$ iteration and $\alpha_n$ controls the enhancement strength.
The process progressively brightens dark regions while limiting
over-exposure in already illuminated regions.

The Zero-DCE network uses the lightweight fully convolutional DCE-Net~\cite{zerodce}, which is a lightweight seven-layer fully convolutional
network operating at input resolution (no pooling/upsampling), enabling
real-time inference ($>$500 FPS) and predicting eight tanh-constrained
RGB curve parameter sets over eight enhancement stages.

\subsection{Object Detection, Depth and Tracking}

\subsubsection{Object Detection using YOLO11n}
Frames satisfying the motion-triage criterion are processed using
YOLO11n \cite{yolov11} to localize objects relevant to subsequent
interaction analysis. The detector produces bounding boxes together
with the corresponding class labels and confidence scores. GPU
acceleration is used when available, with CPU execution provided as a
fallback.

For each detected object, the bounding-box geometry is represented
using its corner coordinates $(x_1,y_1,x_2,y_2)$, from which the
centroid $(c_x,c_y)=\left(\frac{x_1+x_2}{2},\frac{y_1+y_2}{2}\right)$,
width $w=x_2-x_1$, height $h=y_2-y_1$, and area $A=wh$ are derived.
These spatial attributes provide the geometric representation required
by the subsequent stages.

\subsubsection{Semantic Filtering}

A semantic filtering layer is applied to the detections from the object detector in an effort to limit downstream processing to objects relevant to chain-snatching analysis. The filter only keeps the object categories that are relevant for crime: person, motorcycle, bicycle, car, bus and truck. This constraint cuts down irrelevant background objects and environmental clutter irrelevant to the subsequent interaction analysis.

Apart from the semantic class constraint, the detections are required to satisfy a confidence threshold above 0.35 and a minimum bounding-box area criterion; otherwise they are discarded.

\subsubsection{Monocular Depth Estimation using Depth-Anything V2}
After object detection, the proposed framework uses Depth-Anything V2
\cite{yang2024depth} to estimate a dense relative depth map for each
illumination-normalized surveillance frame, using a DINOv2 ViT backbone
for feature extraction and a DPT decoder to generate high-resolution
depth maps from a single RGB image.

The estimated depth map complements the 2D object detections from
YOLOv11 by providing 3D scene context, helping distinguish genuine
physical interactions from apparent intersections caused by perspective
projection. Depth estimation runs in parallel with object detection so
that spatial localization and relative depth are jointly available for
interaction validation; candidate suspect--victim pairs satisfying the
spatial-proximity and depth-consistency criteria are then forwarded for
further analysis.
\subsubsection{Multi-Object Tracking}

Temporal correspondence between detections across consecutive
frames is established by a multi-object tracking layer. The tracker
associates bounding boxes over time and assigns persistent track identities
to individual entities, which allows their trajectories to be maintained across
successive observations and short-term occlusions.

The tracking mechanism combines Kalman-filter-based state prediction with
intersection-over-union (IoU) based data association, following the matching
principles used in modern tracking approaches such as ByteTrack~\cite{bytetrack} and
BoT-SORT. Each active track maintains a persistent identity together with
its historical centroid and bounding-box observations. These trajectories
provide the temporal information required by the subsequent motion-feature
and relationship-analysis stages.

\subsubsection{Motion Feature Extraction}

The temporal trajectories generated by the tracking stage are further
processed to derive motion descriptors for each tracked entity. The
centroid positions of a track across consecutive frames are used to
estimate its instantaneous velocity, average speed over a sliding temporal
window, acceleration, total displacement, and movement direction.

The resulting kinematic descriptors are associated with the corresponding
tracks and provide quantitative information about the movement of each
entity. These features are subsequently used by the relationship and
interaction-analysis stages to characterize relative motion between
tracked entities.

\subsection{Interaction and Behaviour Analysis}

\subsubsection{Relationship Analysis}

To characterize interactions between tracked entities, the framework
performs pairwise spatial and motion analysis between relevant person and
vehicle tracks. For each candidate pair, the Euclidean distance between
their image-space centroids is calculated as

\begin{equation}
d(p_1,p_2)=
\sqrt{(x_2-x_1)^2+(y_2-y_1)^2}.
\end{equation}

The framework additionally evaluates the relative closing speed and
spatial orientation between the entities. A pair is considered for
relationship analysis when its centroid distance falls within a
150-pixel threshold. These measurements provide a representation of the
spatial and relative-motion relationship between tracked entities and
are subsequently used by the interaction-management stage to identify
persistent interactions.

\subsubsection{Interaction Management}
Relationship observations from the preceding stage are maintained over
time using a finite-state interaction manager that tracks person--vehicle
pairs across consecutive frames, distinguishing temporary separations or
occlusions from true interaction termination.

An interaction is initialized in the \textit{NEW} state when a
person--vehicle pair satisfies the spatial proximity criterion, and
transitions to \textit{ACTIVE} after a second consecutive observation. A
temporarily disappearing pair enters the \textit{LINGERING} state,
allowing the relationship to persist through short-term observation gaps;
after 30 inactive frames it transitions to \textit{ENDED} and subsequently
\textit{ARCHIVED}. A lingering interaction returns to \textit{ACTIVE} if
the pair reappears.

For each active interaction, temporal, spatial, and motion histories are
maintained, including distances, centroid positions, speeds, and movement
headings, along with relative kinematic quantities such as closing
velocity, acceleration, heading difference, and trajectory cosine
similarity:
\begin{equation}
S_{\mathrm{traj}} =
\frac{\mathbf{d}_{p}\cdot\mathbf{d}_{v}}
{\|\mathbf{d}_{p}\|\|\mathbf{d}_{v}\|},
\end{equation}
where $\mathbf{d}_{p}$ and $\mathbf{d}_{v}$ denote the displacement
vectors of the person and vehicle, respectively.

A normalized interaction confidence score $C\in[0,1]$ is computed from
proximity, duration, and relative closing speed:
\begin{equation}
C =
\min\left(1.0,\,
0.5S_{\mathrm{proximity}}
+0.3S_{\mathrm{duration}}
+0.2S_{\mathrm{velocity}}\right),
\end{equation}
where the proximity term increases as distance decreases, the duration
term saturates at 60 frames, and the velocity term increases with closing
speed. These histories and confidence estimates are retained as
structured evidence for subsequent stages.
\subsubsection{Behaviour Intelligence Layer}
Interaction histories from the preceding stages are converted into
discrete behavioural primitives using a rule-based behaviour intelligence
layer. At each frame, active interactions are evaluated using spatial
proximity, relative velocity, acceleration, heading, trajectory
similarity, and motion changes; each satisfied rule generates a
time-stamped primitive, representing the interaction's evolution as a
chronological sequence rather than isolated frame-level observations.

The behavioural rules capture phases such as approaching, moving away,
following, parallel motion, close and stationary interaction, rapid
acceleration/deceleration, sudden direction changes, and rapid
separation, using thresholds derived empirically: approaching
(relative velocity $<-2.0$ px/frame), close interaction (distance
$<80$ px), following (trajectory cosine similarity $>0.70$), rapid
acceleration ($>3.0$ px/frame$^2$), and rapid separation (relative
velocity $>5.0$ px/frame).

The resulting primitives are organized chronologically into a behaviour
timeline, subsequently used by the behaviour graph and multimodal fusion
stages to derive higher-level behavioural evidence~\cite{yan2018stgcn}.

\subsubsection{Behaviour Graph Modeling}
The behavioural primitives are organized into a spatio-temporal directed
graph representing the evolution of an interaction, constructed
incrementally by evaluating the active interaction and selecting the
highest-priority behavioural pattern at each frame.

Each node represents a behavioural pattern, storing its type, temporal
extent, confidence, priority, supporting primitives, and associated
motion/spatial evidence; a new node is created when the primary pattern
changes, otherwise the existing node is updated with the latest end
frame, duration, confidence, and kinematic statistics, preserving both
persistent states and transitions between phases. Directed temporal
edges are created on pattern transitions, recording the source and
target patterns, transition frame, timestamp, and transition condition. 

Pattern detection combines spatial, kinematic, alignment, and temporal
criteria—proximity, approach/separation velocity, acceleration,
stationary speed, heading alignment, trajectory similarity, and minimum
duration—drawing on evidence accumulated by the preceding interaction
stages. The resulting behaviour graph constitutes \textit{Stream A} of the proposed two-stream framework

\subsubsection{ROI Selection and Preparation}
To provide stable spatial inputs for pose estimation, regions of
interest (ROIs) corresponding to the person in an active interaction
are extracted from the tracked person's frame-wise bounding boxes,
restricting pose analysis to the relevant participant rather than
unrelated subjects in the scene.

Missing bounding-box observations from temporary occlusions are handled
via linear interpolation for gaps up to 15 frames. Unphysical box-area
variations—deviating from the median area by more than three standard
deviations—are treated as outliers and replaced with the previous valid
bounding box. Temporal jitter is then reduced using Exponential Moving
Average (EMA) smoothing: for a bounding-box vector
$B_t=[x_1,y_1,x_2,y_2]^T$,
\begin{equation}
B_t^{\mathrm{smooth}}
=
\alpha B_t^{\mathrm{raw}}
+
(1-\alpha)B_{t-1}^{\mathrm{smooth}},
\end{equation}
where $\alpha=0.4$.

To retain contextual information such as extended body parts and nearby
vehicle-contact regions, smoothed boxes are expanded by 25\% in width
and height. Resulting ROI sequences are evaluated on sequence
completeness, missing-frame ratio, bounding-box stability, and temporal
coverage, and accepted only when meeting quality criteria, removing
unreliable crops before pose estimation. Accepted ROIs are extracted
and organized into interaction-specific sequences (typically 30--60
frames), yielding interpolated, outlier-filtered, smoothed, expanded,
and quality-validated sequences for the pose-estimation stage.

\subsubsection{Skeleton Sequence Construction}
Temporally smoothed pose keypoints are transformed into normalized
skeleton sequences for graph-based action recognition, using
hip-centred normalisation to reduce the effect of variations in image
position and subject location.

The normalized keypoints are organized into fixed-length sequences of
$T=30$ frames, following the NCTVM tensor convention ($N$: batch,
$C$: coordinate channels, $T$: temporal dimension, $V$: skeletal
joints, $M$: number of persons). For the proposed implementation, the
resulting tensor is
\begin{equation}
(1,C,T,V,1)=(1,2,30,17,1),
\end{equation}
with two coordinate channels (2D joints), 17 joints per skeleton, and a
single subject per sequence. This provides a consistent spatio-temporal
representation for the subsequent ST-GCN action recognition stage.

\subsubsection{ST-GCN Action Recognition}
Normalized skeleton sequences are processed using a Spatial-Temporal
Graph Convolutional Network (ST-GCN)~\cite{yan2018stgcn} to recognize actions
from joint-level motion patterns. The human body is represented as a
graph with the 17 skeletal joints as nodes and anatomical connections
as spatial edges, extended into a spatio-temporal graph via temporal
connections between corresponding joints across consecutive frames.

Spatial graph convolution operates on the normalized adjacency matrix
of the COCO-17 topology, while temporal convolution models joint
configuration evolution across the $T=30$ frame sequence. Each
graph-convolution block is followed by batch normalization and ReLU
activation, with global average pooling applied before a fully
connected layer produces probabilities over ten action classes:
\textit{Walking}, \textit{Standing}, \textit{Running},
\textit{Approaching}, \textit{Reaching}, \textit{Grabbing},
\textit{Pulling}, \textit{Turning}, \textit{Falling}, and
\textit{Unknown}—providing domain-specific motion representations for
distinguishing ordinary pedestrian activity from chain-snatching-related
motion patterns.

The resulting action predictions constitute \textit{Stream B} of the
behavioural analysis framework, providing learned skeleton-based action
evidence in contrast to the rule-based Stream A behaviour graph; the two
streams are subsequently combined by the behaviour fusion stage. Each
inference produces the predicted action, confidence, class
probabilities, top-$k$ predictions, and inference latency, together
with the associated sequence, interaction, and track identifiers,
enabling incorporation into the subsequent behavioural fusion stage.


\subsection{Crime-Specific Snatch Signature Reasoning}
\subsubsection{Snatch Signature Matching}
The fused behavioural evidence is evaluated by a crime-specific
signature engine to determine whether an interaction is consistent with
a chain-snatching event. Rather than directly classifying an entire
video sequence as a crime, the engine evaluates the observed interaction
against predefined behavioural signatures for different snatching
scenarios—separately for motorcycle-based and pedestrian-based cases.
Each signature evaluates the available behavioural evidence using
weighted criteria derived from interaction dynamics and behavioural
patterns, producing a signature score normalized to $[0,1]$ that
represents the degree of agreement with the corresponding pattern.

A score of $S \geq 0.70$ is treated as a high-confidence match. Along
with the numerical score, the engine generates an evidence-based
explanation indicating which signature conditions were satisfied or
not, providing an interpretable decision layer between multimodal
behavioural evidence and the final forensic event record.

\subsubsection{Forensic Event Indexing and Retrieval}
The final stage converts crime-specific assessments from the snatch
signature engine into structured forensic event records, building
sub-millisecond ($<$0.1~ms) $O(1)$ inverted search indices. Each event
retains the decision and associated behavioural and interaction
evidence, enabling detected incidents to be searched and retrieved
without reprocessing the original video data.

A multi-attribute inverted indexing scheme organizes forensic events by
attributes including event and video identifiers, decision, track
identities, behavioural patterns, detected actions, and descriptive
tags, enabling direct retrieval of events matching specific forensic
queries. In addition to the searchable index, the framework preserves
associated evidence artefacts—structured event records, thumbnails, and
annotated video clips—providing investigators with both searchable
metadata and the corresponding visual evidence for forensic examination.

\section{Experimental Results and Discussion}

\subsection{Experimental Setup}

\subsubsection{Dataset}

The framework was evaluated on 43 real-world CCTV video sequences from the Snatch 1.0 benchmark and supplementary surveillance footage, comprising 96,926 frames across 3,877.05~s (64.62~min). The dataset contains 35 positive chain-snatching sequences and 8 negative control sequences of normal surveillance. The recordings encompass heterogeneous surveillance conditions, including varied spatial resolutions, frame rates, static outdoor camera views, daytime, mixed-lighting, and night-time lighting environments.

\subsubsection{Ground Truth and Evaluation Protocol}

Binary video-level ground-truth annotations were provided for all 43 sequences (35 positive, 8 negative), with temporal interaction boundaries established for snatching events. Annotations also cover standard action categories (Walking, Standing, Approaching, Reaching, Grabbing, Pulling, Escaping) and interaction classes. Evaluation measures event-level detection effectiveness (precision, recall, F1-score, accuracy, FPR, FNR), pipeline efficiency (latency, throughput, search-space reduction), stage-wise evidence preservation, and architectural ablation across component configurations.

\subsubsection{Computational Environment}

Experiments were conducted on a CPU--GPU system featuring an 8-core processor and NVIDIA CUDA acceleration. Operational measurements recorded an average CPU utilization of 22.4\%, resident memory of 824.5~MB, GPU VRAM allocation of $\approx$850~MB, and average GPU utilization of 18.5\%.

\subsection{Motion Detection Performance Evaluation}

The keyframe-selection capability of five motion detectors---Baseline (No Filtering), Frame Difference, Gaussian Mixture Model (GMM), K-Nearest Neighbour (KNN), and Mixture of Gaussians Version 2 (MOG2)---was evaluated across 215 benchmark runs (43 videos $\times$ 5 algorithms). Evaluation focuses on processing throughput (FPS), frame reduction (\%), and motion continuity.

\subsubsection{Motion Detection Benchmarking}
For each video, foreground masks identified motion regions to measure
speed (FPS), frame reduction percentage, and motion sequence continuity
required for downstream analysis.

\begin{table}[htbp]
\centering
\scriptsize
\setlength{\tabcolsep}{4pt}
\renewcommand{\arraystretch}{0.9}
\caption{Overall Motion Detection Benchmark Results}
\label{tab:benchmark}
\begin{tabular}{lccc}
\hline
\textbf{Detector} & \textbf{Reduction (\%)} & \textbf{FPS} & \textbf{Continuity} \\
\hline
Baseline          & 0.00  & 1001.35 & 0.995 \\
Frame Difference  & 50.83 & 444.82  & 0.610 \\
GMM               & 34.73 & 114.16  & 0.851 \\
KNN               & 4.00  & 90.77   & 0.992 \\
MOG2              & 5.17  & 144.18  & 0.985 \\
\hline
\end{tabular}
\end{table}
As shown in Table~\ref{tab:benchmark}, Frame
Difference achieves high throughput (444.82 FPS) and reduction
(50.83\%), but its low continuity score (0.610) severely fragmentates
temporal sequences, whereas KNN and MOG2 preserve motion continuity
(0.992 and 0.985, respectively). MOG2 provides the optimal balance of
throughput (144.18 FPS), frame reduction (5.17\%), and continuity
(0.985) without compromising downstream tracking, establishing its
selection for initial motion triage.

A further sensitivity analysis across the motion-area threshold
(0\%--10\%, Fig.~\ref{fig:threshold})—where higher thresholds discard
more static frames but risk fragmenting subtle movements—confirms that
MOG2 maintains stable continuity across moderate threshold settings
while steadily reducing uninformative frames, reinforcing its selection
for the framework.
\begin{figure}[!t]
\centering
\includegraphics[width=0.92\textwidth]{motion_area_threshold.png}
\caption{Motion Area Threshold Sensitivity Analysis}
\label{fig:threshold}
\end{figure}

\subsection{Pipeline Bottleneck and Latency Analysis}

Computational cost was profiled across all pipeline stages, measuring total execution time and per-frame latency (Table~\ref{tab:pipeline_latency} and Fig.~\ref{fig:runtime}).

\begin{table}[t]
\centering
\caption{Computational runtime and latency of major pipeline stages.}
\label{tab:pipeline_latency}
\scriptsize
\setlength{\tabcolsep}{2pt}
\renewcommand{\arraystretch}{0.9}
\begin{tabular}{lrrrr}
\hline
\textbf{Stage} &
\textbf{Time (s)} &
\textbf{Run (\%)} &
\textbf{Avg. (ms)} &
\textbf{Peak (ms)} \\
\hline
Motion & 31.1461 & 10.90 & 13.99 & 72.00 \\
YOLO & 118.7951 & 41.57 & 54.82 & 1943.19 \\
Tracking & 135.3338 & 47.36 & 68.21 & 152.97 \\
Relationship & 0.4648 & 0.16 & 0.23 & 22.00 \\
Events & 0.0000 & 0.00 & 0.00 & 0.00 \\
\hline
\end{tabular}
\end{table}

\begin{figure}[!t]
    \centering
    \includegraphics[width=0.70\columnwidth]{stage_runtime_contribution.png}
    \caption{Computational runtime contribution of the major pipeline stages. Tracking and YOLO detection account for 88.93\% of the measured runtime.}
    \label{fig:runtime}
\end{figure}

The profiling identifies tracking (47.36\% runtime, 68.21~ms avg. latency) and YOLO object detection (41.57\% runtime, 54.82~ms avg. latency) as the primary computational bottlenecks, collectively representing 88.93\% of overall execution time. Motion filtering contributes 10.90\% (13.99~ms avg.), whereas relationship analysis consumes only 0.16\% (0.23~ms avg.). This confirms that spatial relationship reasoning introduces minimal latency while effectively pruning search space prior to complex event synthesis.

\subsection{Pipeline Architecture Ablation}

An ablation study evaluated the progressive integration of system modules across four configurations: (A) YOLO only, (B) Motion + YOLO, (C) Motion + YOLO + Tracking, and (D) Full Pipeline (Table~\ref{tab:ablation}).

\begin{table}[t]
\centering
\caption{Ablation study of the proposed pipeline architecture.}
\label{tab:ablation}
\resizebox{\textwidth}{!}{%
\begin{tabular}{lrrrrr}
\hline
\textbf{Configuration} &
\textbf{Runtime (s)} &
\textbf{Retained} &
\textbf{Detections} &
\textbf{Tracks} &
\textbf{Events} \\
\hline
YOLO Only & 36.51 & 1021 & 2569 & 0 & 0 \\
Motion + YOLO & 95.31 & 1984 & 4947 & 0 & 0 \\
Motion + YOLO + Tracking & 196.0 & 1984 & 4947 & 4022 & 0 \\
Full Pipeline & 187.4 & 53 & 4947 & 4022 & 67 \\
\hline
\end{tabular}%
}
\end{table}

To make un-triaged baseline evaluation computationally practical, Configuration A samples alternate frames ($\texttt{frame\_num}\bmod 2=0$), processing 1,113 of 2,227 frames and yielding 1,021 detection frames. Configurations B and C retain 1,984 candidate frames. The Full Pipeline (Config D) reduces retained candidate frames to 53---a 97.3\% search-space reduction relative to Config C. Crucially, Config D synthesizes 67 structured forensic events (vs. 0 for Configs A--C) with a total runtime of 187.4~s, proving that downstream stages successfully structure raw object tracks into forensic evidence without incurring added latency overhead.

\subsection{End-to-End Forensic Detection Performance}

End-to-end event detection was evaluated across 35 positive and 8 negative sequences (Table~\ref{tab:end_to_end_performance}).

\begin{table}[htbp]
\centering
\caption{End-to-end forensic detection performance.}
\label{tab:end_to_end_performance}
\scriptsize
\renewcommand{\arraystretch}{0.75}
\begin{tabular}{lc}
\hline
\textbf{Metric} & \textbf{Result} \\
\hline
True Positives (TP) & 33 \\
True Negatives (TN) & 8 \\
False Positives (FP) & 0 \\
False Negatives (FN) & 2 \\
Precision & 100.00\% \\
Recall & 94.29\% \\
F1-score & 97.06\% \\
Accuracy & 95.35\% \\
\hline
\end{tabular}
\end{table}

The framework achieved 33 TPs, 8 TNs, 0 FPs, and 2 FNs, corresponding to 100.00\% precision, 94.29\% recall, 97.06\% F1-score, and 95.35\% accuracy. The zero false-positive rate underscores high selectivity during crime-event classification, while missing only two positive events under extreme visual degradation.

\subsection{Progressive Search-Space Reduction}

Search-space filtering was quantified across 12 dataset batches totaling 2,227 frames (Table~\ref{tab:progressive_reduction} and Fig.~\ref{fig:search_space_progression}).

\begin{table}[htbp]
\centering
\scriptsize
\setlength{\tabcolsep}{4pt}
\renewcommand{\arraystretch}{0.85}
\caption{Progressive search-space reduction across validated pipeline stages.}
\label{tab:progressive_reduction}

\begin{tabular}{lcc}
\hline
\textbf{Pipeline Stage} & \textbf{Retained Frames} & \textbf{Search Space (\%)} \\
\hline
Raw Video Input   & 2227 & 100.00 \\
Motion Filtering  & 2167 & 97.31 \\
YOLO Detection    & 1984 & 89.09 \\
Tracking Stage    & 1984 & 89.09 \\
Relationship Engine & 53 & 2.38 \\
Candidate Events  & 53 & 2.38 \\
\hline
\end{tabular}
\end{table}

\begin{figure}[t]
    \centering
    \includegraphics[width=0.70\columnwidth]{search_space_progression.png}
    \caption{Progressive reduction of the CCTV forensic search space across the proposed pipeline stages.}
    \label{fig:search_space_progression}
\end{figure}

The search space decreases from 2,227 initial frames to 2,167 after motion triage (97.31\%) and 1,984 following YOLO detection and tracking (89.09\%). The relationship analysis stage delivers the dominant reduction, pruning candidate frames to 53 (2.38\% of original volume, representing a 97.3\% relative reduction from post-detection). This confirms the architecture's capacity to isolate high-probability temporal windows for downstream reasoning.

\subsection{Evidence Preservation Across Pipeline Stages}

To ensure aggressive filtering does not discard critical evidence, stage-wise preservation was benchmarked across all 35 ground-truth events (Table~\ref{tab:event_preservation}).

\begin{table}[htbp]
\centering
\scriptsize
\setlength{\tabcolsep}{2.5pt}
\renewcommand{\arraystretch}{0.85}
\caption{Stage-wise preservation of ground-truth chain-snatching events.}
\label{tab:event_preservation}

\begin{tabular}{lccccc}
\hline
\textbf{Stage} &
\textbf{GT} &
\textbf{Entering} &
\textbf{Preserved} &
\textbf{Disc.} &
\textbf{Recall (\%)} \\
\hline
Raw Video Input       & 35 & 35 & 0  & 0 & 100.00 \\
Motion Triage         & 35 & 35 & 0  & 0 & 100.00 \\
Semantic Filtering    & 35 & 35 & 0  & 0 & 100.00 \\
Interaction Manager   & 35 & 35 & 0  & 0 & 100.00 \\
Behaviour Graph       & 35 & 35 & 0  & 0 & 100.00 \\
ROI Selection         & 35 & 34 & 0  & 1 & 97.14 \\
Pose Estimation       & 34 & 34 & 0  & 0 & 100.00 \\
Behaviour Fusion      & 34 & 33 & 0  & 1 & 97.06 \\
Signature Engine      & 33 & 33 & 0  & 0 & 100.00 \\
\hline
\end{tabular}
\end{table}

Early filtering stages (Motion Triage through Behaviour Graph) maintained 100.0\% event recall. Initial evidence loss occurred at ROI Selection (34/35 preserved, 97.14\% stage recall), with a second loss at Behaviour Fusion (33/34 preserved, 97.06\% stage recall). Overall cumulative recall reached 94.29\% (33/35 events preserved; 5.71\% loss). Missing events were isolated to severe vehicle occlusion and extreme camera distance rather than early motion pruning.

\subsection{Forensic Indexing and Investigator Retrieval Performance}

The forensic indexing layer was benchmarked across 1,000 queries per category to measure investigator search latency and throughput (Table~\ref{tab:retrieval_performance}).

\begin{table}[htbp]
\centering
\caption{Forensic indexing and investigator retrieval performance.}
\label{tab:retrieval_performance}
\resizebox{\textwidth}{!}{%
\begin{tabular}{lccccc}
\hline
\textbf{Query Type} &
\textbf{Avg. Latency (ms)} &
\textbf{Max. Latency (ms)} &
\textbf{Accuracy (\%)} &
\textbf{Completeness (\%)} \\
\hline
Single Keyword Lookup & 0.042 & 0.12 & 100.0 & 100.0 \\
Multi-Attribute Filter & 0.068 & 0.18 & 100.0 & 100.0 \\
Track ID Trace Lookup & 0.025 & 0.08 & 100.0 & 100.0 \\
Full-Text Boolean Query  & 0.095 & 0.25 & 100.0 & 100.0 \\
\hline
\end{tabular}%
}
\end{table}

Average latency ranged from 0.025~ms (Track ID Trace) to 0.095~ms (Full-Text Boolean), with peak latencies $\le$0.25~ms. All query types achieved 100.0\% accuracy and retrieval completeness, sustaining throughput $>10,000$ searches/sec with a compact memory footprint ($\approx 1.8$~KB per event).

\subsection{Inferential Statistical Significance Validation}

Paired $t$-tests ($N=43$ video streams, $df=42$) evaluated the statistical significance of F1-score improvements across architectural configurations (Table~\ref{tab:statistical_significance}).

\begin{table}[htbp]
\centering
\caption{Statistical significance of progressive architectural improvements in F1-score.}
\label{tab:statistical_significance}
\resizebox{\textwidth}{!}{%
\begin{tabular}{lccccc}
\hline
\textbf{Configuration} &
\textbf{Mean F1} &
\textbf{95\% CI} &
\textbf{$t$-statistic} &
\textbf{$p$-value} &
\textbf{Cohen's $d$} \\
\hline
Config A (Baseline) & 0.6737 & 0.6581--0.6893 & 28.45 & $<0.0001$ & 5.12 \\
Config B (+ Motion) & 0.7347 & 0.7203--0.7491 & 22.14 & $<0.0001$ & 4.35 \\
Config C (+ Graph) & 0.8449 & 0.8332--0.8566 & 14.82 & $<0.0001$ & 2.91 \\
Config D (Proposed) & \textbf{0.9706} & \textbf{0.9643--0.9769} & -- & -- & -- \\
\hline
\end{tabular}%
}
\end{table}

Configuration D achieved the highest mean F1-score (0.9706, 95\% CI: 0.9643--0.9769). Pairwise comparisons against Baselines A, B, and C yielded $p < 0.0001$ in all cases with large effect sizes ($d \ge 2.91$), validating that performance gains stem from architectural design rather than sample variance.

\subsection{Overall Discussion}

The evaluation confirms that progressive multi-stage filtering dramatically reduces forensic search space while preserving crime evidence. Achieving 100.00\% precision and 94.29\% recall across challenging CCTV streams, the framework prunes 97.62\% of raw frames down to 2.38\% candidate frames without sacrificing ground-truth events during early triage. Computational profiling highlights object detection and tracking as the primary runtime contributors (88.93\%), whereas spatial relationship analysis adds negligible overhead (0.16\%) while yielding the highest candidate reduction.

Furthermore, the structured forensic index enables sub-millisecond retrieval across multi-attribute queries, bridging real-time video analytics and forensic investigation. Primary performance limits are tied to extreme visual occlusion and distance during downstream pose/ROI extraction rather than early motion filtering, demonstrating the robust evidence preservation of the proposed architecture.
%%%%%%%%%%%%%%%%%%%%%%%%%%%%%%%%%%%%%%%%%%%%%%%%%%%%%%%%%%%%%

\section{Conclusion}
This work presented an AI-based forensic search framework for automated
analysis of CCTV footage containing chain-snatching events. The proposed
multi-stage architecture progressively processes surveillance footage
through motion filtering, object detection, tracking, spatial
relationship analysis, behavioural reasoning, pose-based action
recognition, evidence fusion, crime-specific signature matching, and
forensic indexing. The framework is designed to reduce computational
workload while preserving evidence relevant to forensic investigation.

The complete system achieved 100.00\% precision, 94.29\% recall,
97.06\% F1-score, and 95.35\% accuracy on the evaluated dataset. The
progressive pipeline reduced the retained search space to 2.38\% of the
original video frames while preserving 33 of 35 ground-truth events.
The results further show that the largest reduction in search space
occurs during relationship analysis rather than the initial filtering
stages, highlighting the importance of higher-level spatial and
behavioural reasoning for narrowing forensic candidates. The forensic
retrieval layer additionally achieved 100.0\% search accuracy and
completeness with sub-millisecond average query latency. These results
demonstrate the feasibility of combining progressive visual and
behavioural analysis with structured forensic retrieval to support
efficient CCTV-based chain-snatching investigation.

\bibliographystyle{splncs04}
\bibliography{ref}

\end{document}